% Chapter 3 -- Building Turbo Initech and its programs on the development
% system.  Sources (Law 1): the Makefile recipes of test-tps-lex-build
% (INITECHC compiles os/tps/tps.pas; links bare and InitechDOS images;
% TPS.COM; the BSS reserve and the 0x6f000 ceiling), test-tps-gen-fpc
% (Free Pascal build of tps.pas; the corpus; self-generation twice; nasm
% -f elf32, ld -m elf_i386 with seed/rt/seed.ld or seed_dos.ld, objcopy -O
% binary), test-tps-gen-os (TPS.COM on InitechDOS).  The make runs printed
% are replayed by test-pubs-turbo; harness/pubs/turbo_replay.sh runs the same
% assemble/link commands for every program printed in the manuals.
\Gates{test-tps-lex-build, test-tps-gen-fpc, test-tps-gen-build, test-tps-gen-os, test-pubs-turbo}
\chapter{Building Turbo Initech}
\ix{building Turbo Initech}

All the steps of this chapter run on the development system, from the
directory that holds the InitechOS sources, with the \code{make} program
and the targets named here.

\section{The pieces}
\begin{BorTable}{Sources and products}{p{36mm}p{62mm}}
\Hd{Piece} & \Hd{What it is}\EndHd
\file{seed/} & the C source of the genesis compiler \file{INITECHC}\\
\file{os/tps/tps.pas} & the Pascal source of Turbo Initech, one file\\
\file{seed/rt/} & the run-time library: the InitechDOS start object, the
  stand-alone start object, the file object, and the two link maps\\
\file{os/tps/fixtures/} & the corpus of test programs (Appendix A)\\
\file{TPS.COM} & the Turbo Initech compiler for InitechDOS, made by
  \file{INITECHC}: $K_1$\\
\end{BorTable}

\section{Making the first compiler}
\ix{K1}
\file{INITECHC} compiles \file{tps.pas} to assembly language, which is
assembled and linked twice: once with the stand-alone start object, and
once with the InitechDOS start object and the file object to make
\file{TPS.COM}.  The build checks that the compiler's global variables stay
within their reserve and that the whole image ends below 6F000h, the top of
the InitechDOS program area:

\MakeRun{tps-lex-build}

\section{Compiling a program on the development system}
\ix{assembling}
\ix{linking}
The development system also runs Turbo Initech itself, in a build made by
the Free Pascal compiler from the same \file{tps.pas}; the InitechDOS
compiler \file{TPS.COM} is checked to print the same output.  A program is
compiled and made into a \file{.COM} file with these commands, which are
the ones the build runs:

\begin{BorTable}{From TPSIN.PAS to a .COM file}{p{20mm}p{78mm}}
\Hd{Step} & \Hd{Command}\EndHd
Compile & run Turbo Initech in the directory that holds
  \file{TPSIN.PAS}\\
Assemble & \code{nasm -f elf32 TPSOUT.S -o TPSOUT.o}\\
Link & \texttt{ld -m elf\_i386 -T seed/rt/seed\_dos.ld -o PROG.ELF
  build/seed\_rt\_dos.o build/seed\_fileio\_rt.o TPSOUT.o}\\
Convert & \code{objcopy -O binary PROG.ELF PROG.COM}\\
\end{BorTable}

The InitechDOS start object must be first.  To test a program without
InitechDOS, link it instead with \file{seed/rt/seed.ld} and the stand-alone
start object \file{build/seed\_rt.o}; a program that uses files needs
InitechDOS.  Every program printed in the Turbo Initech manuals is built
with these commands, run, and compared with its printed output each time
the manuals are checked.

\section{Turbo Initech compiling its own source}
\ix{self-generation}
The build compiles every program of the corpus, and then the Turbo
Initech source itself, twice each, and requires the two results to be the
same bytes:

\MakeRun{tps-gen-fpc}

The last lines give the size of the assembly language Turbo Initech writes
for its own source and show that the linked compiler ends below 6F000h.
